첫 번째 줄에 계획의 개수 $N$과 질의의 개수 $Q$가 주어진다. $(1 \le N, Q \le 10\,080)$

이후 $N$개의 줄에 기룡이의 계획이 아래 형식으로 주어진다.
$$\texttt{DAY HH:MM HH:MM SUBJECT}$$

그 이후 $Q$개의 줄에 확인하고 싶은 시점이 아래 형식으로 주어진다.
$$\texttt{DAY HH:MM}$$

\tt{DAY}는 \tt{MON}, \tt{TUE}, \tt{WED}, \tt{THU}, \tt{FRI}, \tt{SAT}, \tt{SUN} 중 하나이며, 모든 시각은 24시간 형식을 사용한다. $(00 \le \texttt{HH} \le 23, 00 \le \texttt{MM} \le 59)$

\tt{SUBJECT}는 길이가 1 이상 30 이하의 알파벳 소문자와 밑줄(\_)로만 이루어진 문자열이다.

모든 계획은 월요일부터 일요일까지 시간 순으로 주어진다. 한 계획 내에서 시작 시각은 종료 시각보다 늦지 않고, 하나의 계획이 자정을 넘어가는 경우는 없다. 서로 다른 두 계획의 진행 시간이 겹치는 경우는 주어지지 않는다.